\documentclass[12pt,a4paper]{article}

\usepackage{template}
\usepackage{amsmath}
\usepackage{amssymb}
\usepackage{booktabs}

\newcommand{\ord}{\operatorname{ord}}

\hypersetup{
  colorlinks=true,
  linkcolor=blue!45!black,
  urlcolor=blue!45!black,
  citecolor=blue!45!black,
  pdfauthor={Иван Афанасьев},
  pdftitle={Задача 1. Рациональная функция (элементарная версия)}
}

\begin{document}
\sloppy

\begin{center}
  {\LARGE\bfseries Задача 1. Рациональная функция}\\[4pt]
  {\large элементарная версия (только школьные средства)}\\[6pt]
  {\large Иван Афанасьев}\\[2pt]
  {\normalsize Вступительный набор <<group2026famcs>>, ФПМИ БГУ, 2026}
\end{center}

\vspace{6pt}
\begin{center}
\fbox{\begin{minipage}{0.92\textwidth}
\small
\textbf{Использование ИИ.} Решение подготовлено при активном участии
\textbf{GLM-5.2} (помощь в проверке и оформлении гипотез) и
\textbf{OpenCode-CI-Agent} выполнял независимую рецензию на изменения во время Pull Request'а.
\end{minipage}}
\end{center}
\vspace{10pt}

\section{Постановка задачи и договорённости}

Комплексное число $\xi$ называется \emph{корнем $n$-й степени из единицы},
если $\xi^{n}=1$. Рациональная функция над $\mathbb Q$~--- это отношение
$R(x)=P(x)/Q(x)$ многочленов с рациональными коэффициентами, $Q\ne0$;
всегда будем считать дробь сокращённой (у $P$ и $Q$ нет общих многочленных
делителей). Требуется найти все $R$, у которых для любого натурального $n$
и любого корня $\xi$ степени $n$ значение $R(\xi)$ определено и само
является корнем степени $n$:
\[
  \xi^{n}=1 \ \Longrightarrow\ R(\xi)^{n}=1. \tag{$\star$}
\]

\section{Краткий ответ}

\begin{enumerate}[label=\arabic*.]
  \item Корней $n$-й степени из единицы ровно $n$.
  \item Дробно-линейных решений два: $R(x)=x$ и $R(x)=1/x$.
  \item Произвольных рациональных решений~--- в точности одночлены
        $R(x)=x^{m}$, $m\in\mathbb Z$.
  \item[4--5.] Над $\mathbb Q(i)$ и над числовыми полями ответ тот же.
  \item[6.] Над $\mathbb Z_p$ решений становится очень много.
\end{enumerate}

\section{Пункт 1. Корней ровно \texorpdfstring{$n$}{n}}

\emph{Не больше $n$:} по теореме Безу многочлен $x^{n}-1$ степени $n$ имеет
не более $n$ корней.

\emph{Ровно $n$:} числа
$\xi_j=\cos\frac{2\pi j}{n}+i\sin\frac{2\pi j}{n}$, $j=0,1,\dots,n-1$,
попарно различны, и по формуле Муавра $\xi_j^{\,n}=1$.

\section{Пункт 2. Дробно-линейные функции}

Пусть $R(x)=\dfrac{ax+b}{cx+d}$, $ad-bc\ne0$. Используем $(\star)$ для
$\xi=1,-1,i$. Получаем $R(x)=x$ или $R(x)=1/x$.

\section{Пункт 3. Произвольные рациональные функции}

\textbf{Теорема.} Если $R\in\mathbb Q(x)$ удовлетворяет $(\star)$, то
$R(x)=x^{m}$ для некоторого $m\in\mathbb Z$.

\paragraph{Лемма Ф.} Пусть $p$~--- простое и
$\Phi(x)=x^{p-1}+x^{p-2}+\dots+x+1$. Тогда $\Phi$ неприводим над $\mathbb Q$.

\paragraph{Лемма Д (о делимости).} Если $F\in\mathbb Q[x]$ и $F(\zeta)=0$,
где $\zeta$~--- корень степени $p$, то $\Phi\mid F$.

\paragraph{Доказательство теоремы.} Пусть $R=P/Q$~--- сокращённая дробь,
$d_P=\deg P$, $d_Q=\deg Q$. Возьмём простое $p> d_P+d_Q+2$ и корень
$\zeta$ степени $p$. По $(\star)$ имеем $R(\zeta)^{p}=1$, значит
$R(\zeta)=\zeta^{k}$ для некоторого $k$.

Многочлен $F(x)=P(x)-x^{k}Q(x)$ имеет $F(\zeta)=0$. Если $F\equiv0$, то
$R(x)=x^{k}$. Иначе по Лемме~Д $\Phi\mid F$, откуда $p-1\le\deg F$.

Положим $s:=p-k$; тогда $1\le s\le d_Q+1$. По Лемме~Д $P(\xi)=\xi^{k}Q(\xi)$
для всех $\xi=\zeta^{a}$. Умножив на $\xi^{s}$, получим
$\xi^{s}P(\xi)=Q(\xi)$ для $p-1$ различных точек, откуда
$G(x)=x^{s}P(x)-Q(x)\equiv0$ и $R(x)=x^{-s}$.

\section{Пункты 4--5. Коэффициенты из \texorpdfstring{$\mathbb Q(i)$}{Q(i)} и числовых полей}

Одночлены $x^{m}$ по-прежнему подходят. Для $R$ с коэффициентами из
$\mathbb Q(i)$ обозначим через $\overline R$ сопряжённую функцию.
Произведение $S=R\cdot\overline R$ имеет рациональные коэффициенты и
удовлетворяет $(\star)$, поэтому $S=x^{m}$. Отсюда сильные ограничения
на нули и полюсы $R$; полное доказательство $R=x^{k}$ требует средств
выше школьных.

\section{Пункт 6. Конечное поле \texorpdfstring{$\mathbb Z_p$}{Z\_p}}

По малой теореме Ферма $\xi^{p-1}\equiv1$ для $\xi\ne0$. Условие $(\star)$
равносильно $\ord(R(\xi))\mid\ord(\xi)$. Элементов порядка $d\mid p-1$
ровно $d$. Число допустимых функций равно $\prod_{\xi\ne0}\ord(\xi)$,
и каждая реализуется многочленом Лагранжа.

\begin{tabular}{@{}rrr@{}}
\toprule
$p$ & всего решений & из них одночленов \\
\midrule
3  & 2 & 2 \\
5  & 32 & 4 \\
7  & 648 & 6 \\
11 & 12\,500\,000 & 10 \\
13 & 214\,990\,848 & 12 \\
\bottomrule
\end{tabular}

\section{Вычислительная проверка}

Программа \texttt{verify.cpp} перебором подтверждает пункты 2--3 и таблицу
пункта 6.

\section*{Приложение. Неприводимость \texorpdfstring{$\Phi$}{Phi}}

Сдвиг $x=y+1$ даёт $\Phi(y+1)=\sum_{k=1}^{p}\binom{p}{k}y^{k-1}$.
Все коэффициенты, кроме старшего, делятся на $p$, а свободный член
$p$ на $p^{2}$ не делится. По лемме Гаусса и критерию Эйзенштейна $\Phi$
неприводим.

\end{document}
